\documentclass[aspectratio=169]{beamer}
\usetheme{Madrid}
\usecolortheme{seahorse}
\usepackage{graphicx}
\usepackage{booktabs}
\usepackage{array}
\usepackage{amsmath}
\usepackage{xcolor}

\title[Chapter 9: Essential Probability Rules]{\textbf{Chapter 9: Essential Probability Rules}}
\subtitle{School of Mathematical and Computational Sciences \\ University of Prince Edward Island}
\author{Ali Raisolsadat}
\date{Part II: Probability and Distributions}

% Probability notation used throughout the chapter.
\newcommand{\prob}[1]{\operatorname{P}\!\left(#1\right)}
\newcommand{\E}{\operatorname{E}}

% Use the R-generated figure when it is present. This fallback only keeps the
% file compilable when the PNGs have not yet been generated in the same folder.
\newcommand{\lecturefigure}[2][0.72\textheight]{%
  \IfFileExists{#2}{%
    \includegraphics[height=#1,width=\linewidth,keepaspectratio]{#2}%
  }{%
    \fbox{\parbox[c][0.34\textheight][c]{0.82\linewidth}{\centering
      \textbf{Figure file not found}\\[0.2cm]
      {\tiny\texttt{\detokenize{#2}}}\\[0.2cm]
      \small Generate the plot from the Chapter 9 RMarkdown file and place the PNG beside this \texttt{.tex} file.}}%
  }%
}

\begin{document}

% =============================================================
% SLIDE 1: Title
% =============================================================

\begin{frame}
  \titlepage
\end{frame}

%\begin{frame}{Learning Objectives}
%By the end of this chapter, you should be able to:
%\begin{itemize}
%  \item explain randomness and probability using long-run relative frequency;
%  \item distinguish theoretical, empirical, and personal probability;
%  \item identify sample spaces, events, and probability models;
%  \item apply the four fundamental probability rules;
%  \item distinguish discrete and continuous probability models;
%  \item calculate the mean and standard deviation of a discrete random variable;
%  \item convert between probability, risk, and odds.
%\end{itemize}
%\end{frame}

% =============================================================
% PART I: THE IDEA OF PROBABILITY
% =============================================================

\begin{frame}{Randomness and Probability}
\small
When people hear the word ``random,'' they often think of something completely
haphazard or patternless. In statistics, randomness has a more precise meaning.

\vspace{0.2cm}
\begin{block}{Random phenomenon}
A phenomenon is \textbf{random} when individual outcomes are uncertain, but a
regular and predictable distribution emerges over many repetitions.
\end{block}

\begin{block}{Probability}
The probability of an outcome is the proportion of times that outcome would
occur in a very long series of repetitions.
\end{block}

\begin{exampleblock}{Example: Births}
We cannot predict the sex of the next baby with certainty. However, when we
look across a very large number of births, the proportions become remarkably
stable. Probability describes this \emph{long-run regularity}.
\end{exampleblock}
\end{frame}

\begin{frame}{Two Ways to Determine Probability}
How do we decide what probability to assign to an event? In practice, we
usually rely on one of two approaches.

\vspace{0.2cm}
\begin{columns}[T,onlytextwidth]
\column{0.48\textwidth}
\begin{block}{Theoretical probability}
Derived from an underlying mathematical, physical, or biological model.

\medskip
\textbf{Example:} A simplified genetic model suggests
\[
\prob{\text{boy}}=0.50.
\]

The probability comes from the \textbf{model}, not from counting observations.
\end{block}

\column{0.48\textwidth}
\begin{block}{Empirical probability}
Estimated from the long-run frequency observed in real data.

\medskip
\textbf{Example:} Historical birth records may give
\[
\widehat{\prob{\text{boy}}}\approx 0.512.
\]

The probability is estimated from \textbf{data}.
\end{block}
\end{columns}

\vspace{0.2cm}
Both approaches can be used to construct a \textbf{probability model}, and
both must obey the same mathematical probability rules.
\end{frame}

\begin{frame}{The Frequentist Interpretation}
\begin{block}{Core idea}
In the frequentist interpretation, probability is the \textbf{long-run relative
frequency} of an outcome over many repeated trials under similar conditions.
\end{block}

\begin{itemize}
  \item Individual trials remain uncertain.
  \item Small samples can vary substantially from one repetition to another.
  \item As the number of repetitions grows, relative frequencies tend to stabilize.
  \item A probability of $0.5$ does \alert{not} mean exactly half of the next ten trials must be successes.
\end{itemize}

\begin{exampleblock}{Common cold example}
Suppose a survey of $14{,}169$ adults finds that $11.2\%$ had a cold yesterday.
A second random sample would not be expected to give exactly $11.2\%$, but
many repeated samples would reveal a stable population pattern.
\end{exampleblock}
\end{frame}

\begin{frame}{Long-Run Regularity: Coin Toss Simulation}
\begin{columns}[c,onlytextwidth]
\column{0.66\textwidth}
\centering
\lecturefigure[0.62\textheight]{fig_ch9_cointoss.png}

\column{0.32\textwidth}
\small
\begin{itemize}
  \item Each line is a separate sequence of $5{,}000$ coin tosses.
  \item Early proportions can swing sharply because only a few tosses have occurred.
  \item As the number of tosses grows, both proportions move toward $0.5$.
  \item The horizontal axis is logarithmic so the early variation is easier to see.
\end{itemize}
\end{columns}
\end{frame}

\begin{frame}[fragile]{R Example: Simulating Coin Tosses}
The simulation behind the previous figure uses a vector of $0$s and $1$s,
where $1$ represents heads.

\vspace{0.2cm}
\scriptsize
\begin{verbatim}
set.seed(42)
n_tosses <- 5000

# Simulate two independent trials
tosses_A <- sample(c(0, 1), n_tosses, replace = TRUE)
tosses_B <- sample(c(0, 1), n_tosses, replace = TRUE)

# Running proportion of heads after each toss
prop_A <- cumsum(tosses_A) / (1:n_tosses)
prop_B <- cumsum(tosses_B) / (1:n_tosses)
\end{verbatim}

\normalsize
\begin{block}{What the R functions do}
\texttt{sample()} generates the random tosses, while \texttt{cumsum()} keeps a
running total of the number of heads.
\end{block}
\end{frame}

\begin{frame}{The Law of Large Numbers}
\begin{block}{Law of Large Numbers}
As the number of independent repetitions increases, the observed relative
frequency tends to approach the underlying probability.
\end{block}

\begin{itemize}
  \item Early results may fluctuate dramatically because each new observation has a large effect.
  \item With thousands of repetitions, one additional result has very little influence on the cumulative proportion.
  \item This long-run stabilization is one reason large samples are generally more dependable than very small samples.
\end{itemize}

\begin{alertblock}{Important caution}
The Law of Large Numbers describes \textbf{long-run behaviour}. It does not
force a short sequence to ``balance out.'' A coin has no memory of previous tosses.
\end{alertblock}
\end{frame}

\begin{frame}{The Illusion of Randomness}
Which sequence is more likely when a fair coin is tossed six times?

\vspace{0.4cm}
\begin{columns}[T,onlytextwidth]
\column{0.46\textwidth}
\begin{block}{Sequence A}
\centering
\Large H T H T T H
\end{block}

\column{0.46\textwidth}
\begin{block}{Sequence B}
\centering
\Large T T T H H H
\end{block}
\end{columns}

\pause
\vspace{0.4cm}
\begin{alertblock}{Answer}
Both specific sequences have probability
\[
\left(\frac12\right)^6=\frac1{64}.
\]
True randomness naturally produces streaks and clusters.
\end{alertblock}
\end{frame}

\begin{frame}{Personal Probability}
\begin{block}{Definition}
A \textbf{personal probability} is a number between $0$ and $1$ that expresses an individual's scientifically informed judgement about how likely an outcome is.
\end{block}

\begin{exampleblock}{Example}
What is the probability that humans discover extraterrestrial life during this century?
\end{exampleblock}

This event cannot be repeated thousands of times. The probability must therefore be based on evidence, scientific models, and expert judgement rather than direct long-run repetition.

\vspace{0.25cm}
Regardless of its interpretation, every valid probability must obey the mathematical rules introduced later in this chapter.
\end{frame}

% =============================================================
% PART II: PROBABILITY MODELS
% =============================================================

\begin{frame}{Vocabulary of Probability Models}
\small
To use probability mathematically, we need to describe both the possible
outcomes and how likely those outcomes are.

\begin{block}{Sample space}
The \textbf{sample space} $S$ is the set of \emph{all} possible outcomes of a
random phenomenon.
\end{block}

\begin{block}{Event}
An \textbf{event} is one outcome or a collection of outcomes; mathematically,
an event is a subset of $S$.
\end{block}

\begin{block}{Probability model}
A probability model specifies:
\begin{enumerate}
  \item a sample space $S$;
  \item a rule for assigning probabilities to events.
\end{enumerate}
\end{block}

\small
For example, if we record only sex at birth, a simple sample space is
$S=\{\text{Male},\text{Female}\}$.
\end{frame}

\begin{frame}{Example 9.5: Blood Types}
For the eight blood types,
\[
S=\{O+,O-,A+,A-,B+,B-,AB+,AB-\}.
\]

\begin{table}
\centering
\small
\begin{tabular}{c*{8}{>{\centering\arraybackslash}p{0.075\textwidth}}}
\toprule
Type & $O+$ & $O-$ & $A+$ & $A-$ & $B+$ & $B-$ & $AB+$ & $AB-$\\
\midrule
Probability & .39 & .01 & .27 & .005 & .25 & .004 & .07 & .001\\
\bottomrule
\end{tabular}
\end{table}

\begin{itemize}
  \item Each listed outcome is discrete.
  \item The probabilities are nonnegative and sum to $1$.
\end{itemize}
\end{frame}

\begin{frame}{Blood-Type Probability Model}
\begin{columns}[c,onlytextwidth]
\column{0.66\textwidth}
\centering
\lecturefigure[0.62\textheight]{fig_ch9_blood_types.png}

\column{0.32\textwidth}
\small
The bars visualize the probability assigned to each discrete outcome.

\vspace{0.2cm}
\begin{itemize}
  \item $O+$ has the largest probability in this model.
  \item Rare outcomes such as $AB-$ still belong to the sample space.
  \item Adding the eight probabilities gives exactly $1$.
\end{itemize}
\end{columns}
\end{frame}

\begin{frame}[fragile]{R Example: Representing a Discrete Probability Model}
A probability model can be stored directly as vectors or as a data frame.

\vspace{0.15cm}
\scriptsize
\begin{verbatim}
blood_data <- data.frame(
  Type = factor(
    c("O+", "O-", "A+", "A-", "B+", "B-", "AB+", "AB-"),
    levels = c("O+", "O-", "A+", "A-", "B+", "B-", "AB+", "AB-")
  ),
  Probability = c(0.39, 0.01, 0.27, 0.005,
                  0.25, 0.004, 0.07, 0.001)
)

sum(blood_data$Probability)
# 1
\end{verbatim}

\normalsize
\begin{block}{Why check the sum?}
A valid discrete probability model must assign probabilities that sum to $1$.
\end{block}
\end{frame}

\begin{frame}{Changing the Scope of a Model}
Suppose we care only about the Rhesus factor.

\[
S=\{Rh+,Rh-\}.
\]

\begin{table}
\centering
\begin{tabular}{lcc}
\toprule
Rh factor & $Rh+$ & $Rh-$\\
\midrule
Probability & $0.98$ & $0.02$\\
\bottomrule
\end{tabular}
\end{table}

\begin{block}{Key idea}
The sample space depends on how the random phenomenon is defined. A broader model can often be collapsed into a simpler model by combining outcomes.
\end{block}
\end{frame}

\begin{frame}{Example 9.6: A Random Number from 0 to 1}
Let $Y$ be a random number selected from the interval $[0,1]$.

\[
S=\{y:0\le y\le 1\}.
\]

\begin{columns}[c,onlytextwidth]
\column{0.50\textwidth}
\small
\begin{itemize}
  \item There are infinitely many possible decimal values.
  \item We cannot list every possible value and attach a positive probability to each one.
  \item Instead, probabilities are assigned to \textbf{intervals} as areas under a density curve.
  \item The entire area under the density curve must equal $1$.
\end{itemize}

\column{0.48\textwidth}
\centering
\lecturefigure[0.42\textheight]{fig_ch9_continuous_uniform.png}
\end{columns}
\end{frame}

\begin{frame}{Discrete vs. Continuous Sample Spaces}
\begin{columns}[T,onlytextwidth]
\column{0.48\textwidth}
\begin{block}{Discrete}
\begin{itemize}
  \item Finite or countable outcomes
  \item Individual outcomes may have positive probability
  \item Probabilities can be listed in a table
  \item Example: blood type
\end{itemize}
\end{block}

\column{0.48\textwidth}
\begin{block}{Continuous}
\begin{itemize}
  \item An unbroken interval of values
  \item Individual points have probability $0$
  \item Probabilities are areas under a curve
  \item Example: exact height
\end{itemize}
\end{block}
\end{columns}
\end{frame}

% =============================================================
% PART III: PROBABILITY RULES
% =============================================================

\begin{frame}{Four Fundamental Probability Rules}
These rules must hold for \textbf{every} valid probability model, whether the
sample space is discrete or continuous.

\small
\begin{enumerate}
  \item \textbf{Probabilities are proportions:}
  \[
  0\le \prob{A}\le 1.
  \]
  A probability can never be negative or exceed $1$.

  \item \textbf{The entire sample space has probability 1:}
  \[
  \prob{S}=1.
  \]

  \item \textbf{Add disjoint events:}
  \[
  \prob{A\text{ or }B}=\prob{A}+\prob{B}.
  \]

  \item \textbf{Use the complement:}
  \[
  \prob{A^c}=1-\prob{A}.
  \]
\end{enumerate}
\end{frame}

\begin{frame}{Disjoint Events and the Addition Rule}
Two events are \textbf{disjoint} when they have no outcomes in common.

\[
\prob{A\text{ and }B}=0.
\]

Therefore,
\[
\prob{A\text{ or }B}=\prob{A}+\prob{B}.
\]

For several pairwise disjoint events,
\[
\prob{A\text{ or }B\text{ or }C}
=\prob{A}+\prob{B}+\prob{C}.
\]

\begin{alertblock}{Do not double-count}
This simple addition rule applies only when the events cannot occur together.
\end{alertblock}
\end{frame}

\begin{frame}{The Meaning of ``OR'' and ``AND''}
\begin{columns}[T,onlytextwidth]
\column{0.48\textwidth}
\begin{block}{OR}
Inclusive: at least one event occurs.

\medskip
$A$ alone, $B$ alone, or both may occur.
\end{block}

\column{0.48\textwidth}
\begin{block}{AND}
Intersection: both events occur at the same time.

\medskip
This corresponds to the common outcomes of $A$ and $B$.
\end{block}
\end{columns}

\vspace{0.4cm}
When events overlap, the general addition rule is needed to correct for double-counting. That rule is introduced in Chapter 10.
\end{frame}

\begin{frame}{Example 9.7: Positive Rhesus Factor}
The positive blood-type outcomes are disjoint:
\[
O+,\quad A+,\quad B+,\quad AB+.
\]

Therefore,
\begin{align*}
\prob{Rh+}
&=\prob{O+}+\prob{A+}+\prob{B+}+\prob{AB+}\\
&=0.39+0.27+0.25+0.07\\
&=\boxed{0.98}.
\end{align*}
\end{frame}

\begin{frame}{Example 9.7: Complement Rule}
The complement of having a positive Rhesus factor is having a negative Rhesus factor.

\begin{align*}
\prob{Rh-}
&=1-\prob{Rh+}\\
&=1-0.98\\
&=\boxed{0.02}.
\end{align*}

\begin{block}{Why complements are useful}
It is often easier to subtract one probability from $1$ than to add many separate probabilities.
\end{block}
\end{frame}

% =============================================================
% PART IV: DISCRETE AND CONTINUOUS PROBABILITY MODELS
% =============================================================

\begin{frame}{Discrete Probability Models}
\begin{block}{Definition}
A discrete probability model has a finite or countable list of outcomes. The probability of an event is the sum of the probabilities of the outcomes belonging to that event.
\end{block}

A valid discrete probability model must satisfy:
\[
0\le p_i\le 1
\qquad\text{and}\qquad
\sum_i p_i=1.
\]
\end{frame}

\begin{frame}{Example 9.8: Hearing Impairment in Dalmatians}
A study of $5{,}333$ Dalmatians recorded the number of ears with hearing
impairment. Let $X$ be the number of impaired ears for a randomly selected dog.

\begin{columns}[c,onlytextwidth]
\column{0.42\textwidth}
\begin{table}
\centering
\begin{tabular}{lccc}
\toprule
$X$ & $0$ & $1$ & $2$\\
\midrule
$\prob{X=x}$ & $0.70$ & $0.22$ & $0.08$\\
\bottomrule
\end{tabular}
\end{table}

\small
The possible outcomes are the whole numbers $0$, $1$, and $2$, so $X$ is a
\textbf{discrete random variable}. Notice that
\[
0.70+0.22+0.08=1.
\]

\column{0.56\textwidth}
\centering
\lecturefigure[0.48\textheight]{fig_ch9_dalmatian.png}
\end{columns}
\end{frame}

\begin{frame}[fragile]{R Example: A Discrete Probability Distribution}
The Dalmatian probability model can be represented with two short vectors.

\vspace{0.2cm}
\scriptsize
\begin{verbatim}
# Possible values of X
x <- c(0, 1, 2)

# Corresponding probabilities
p_x <- c(0.70, 0.22, 0.08)

# Check that this is a valid probability model
sum(p_x)
# 1

# Probability of at least one impaired ear
sum(p_x[x >= 1])
# 0.30
\end{verbatim}

\normalsize
\begin{block}{Connection to the mathematics}
For a discrete model, the probability of an event is found by
\textbf{adding the probabilities of the outcomes inside the event}.
\end{block}
\end{frame}

\begin{frame}{Some Hearing Impairment}
``Some hearing impairment'' means one or two impaired ears:
\begin{align*}
\prob{X\ge 1}
&=\prob{X=1}+\prob{X=2}\\
&=0.22+0.08\\
&=\boxed{0.30}.
\end{align*}

\begin{alertblock}{Inequality signs matter}
\[
\prob{X\ge 1}=0.30,
\qquad
\prob{X>1}=\prob{X=2}=0.08.
\]
\end{alertblock}
\end{frame}

\begin{frame}{Continuous Probability Models and Density Curves}
A continuous sample space contains an entire interval of possible values, so
we cannot assign probability by listing individual points.

\begin{block}{Density curve}
A density curve:
\begin{enumerate}
  \item lies on or above the horizontal axis;
  \item has total area exactly equal to $1$.
\end{enumerate}
\end{block}

For a continuous random variable $Y$,
\[
\prob{a\le Y\le b}
=\text{area under the density curve from }a\text{ to }b.
\]

\begin{alertblock}{Important}
Probability is represented by \textbf{area}, not by the height of the curve at
one particular value.
\end{alertblock}
\end{frame}

\begin{frame}{Example 9.9--9.10: Women's Heights}
\begin{columns}[c,onlytextwidth]
\column{0.66\textwidth}
\centering
\lecturefigure[0.62\textheight]{fig_ch9_heights.png}

\column{0.32\textwidth}
\small
The histogram shows simulated heights of women aged 40--49, while the smooth
curve provides a continuous probability model.

\vspace{0.2cm}
\begin{itemize}
  \item The total area under the smooth curve is $1$.
  \item The shaded region to the left of $62$ inches represents a probability.
  \item In this model, that area is approximately $0.316$.
\end{itemize}
\end{columns}
\end{frame}

\begin{frame}{Probability as Area}
For the model of women's heights shown previously,
\[
\prob{Y\le 62}=0.316.
\]

\begin{itemize}
  \item The probability is an \textbf{area}, not the height of the curve.
  \item The total area under the curve is $1$.
  \item The shaded area represents $31.6\%$ of the modeled population.
\end{itemize}
\end{frame}

\begin{frame}{Example 9.11: Uniform Density}
For a uniform random number between $0$ and $1$, the density has constant
height $1$. Therefore, probability is especially easy to compute:
\[
\text{area}=\text{width}\times\text{height}.
\]

\begin{columns}[c,onlytextwidth]
\column{0.60\textwidth}
\centering
\lecturefigure[0.52\textheight]{fig_ch9_uniform_shaded.png}

\column{0.37\textwidth}
\begin{align*}
\prob{0.3\le Y\le 0.7}
&=(0.7-0.3)(1)\\
&=\boxed{0.4}.
\end{align*}

\small
The interval occupies $40\%$ of the full width of $[0,1]$, so it also contains
$40\%$ of the total probability.
\end{columns}
\end{frame}

\begin{frame}[fragile]{R Example: Drawing a Uniform Probability Area}
The shaded region in the uniform-density figure is produced by drawing one
rectangle for the full density and another for the probability interval.

\vspace{0.15cm}
\scriptsize
\begin{verbatim}
p_uniform_shaded <- ggplot() +
  geom_rect(
    aes(xmin = 0, xmax = 1, ymin = 0, ymax = 1),
    fill = NA, color = "black"
  ) +
  geom_rect(
    aes(xmin = 0.3, xmax = 0.7, ymin = 0, ymax = 1),
    fill = "#4C72B0", alpha = 0.5
  ) +
  annotate("text", x = 0.5, y = 0.5,
           label = "Area = 0.4")
\end{verbatim}

\normalsize
The code visually connects the probability calculation $0.7-0.3=0.4$ with
the corresponding area under the density curve.
\end{frame}

\begin{frame}{Adding Disjoint Continuous Intervals}
For $Y\sim\operatorname{Uniform}(0,1)$,
\begin{align*}
\prob{Y\le 0.5\text{ or }Y>0.8}
&=\prob{Y\le 0.5}+\prob{Y>0.8}\\
&=0.5+0.2\\
&=\boxed{0.7}.
\end{align*}

The intervals are disjoint, so their areas can be added directly.
\end{frame}

\begin{frame}{The Probability of an Exact Value}
For a continuous random variable,
\[
\prob{Y=c}=0
\]
for every exact value $c$.

\begin{block}{Why?}
A single point has zero width, so the area above that point is zero.
\end{block}

Consequently,
\[
\prob{Y\le 0.8}=\prob{Y<0.8}.
\]
\end{frame}

% =============================================================
% PART V: RANDOM VARIABLES
% =============================================================

\begin{frame}{Random Variables}
Random phenomena often produce numerical outcomes. We use a random variable
to give those outcomes a convenient mathematical name.

\begin{block}{Definition}
A \textbf{random variable} is a variable whose numerical value is determined by
the outcome of a random phenomenon. We usually denote random variables by
uppercase letters such as $X$ or $Y$.
\end{block}

\begin{block}{Probability distribution}
The probability distribution of $X$ specifies:
\begin{itemize}
  \item which values $X$ can take;
  \item how probabilities are assigned to those values.
\end{itemize}
\end{block}

\begin{itemize}
  \item $X=$ number of impaired ears in a Dalmatian: $X\in\{0,1,2\}$.
  \item $Y=$ random number selected from $[0,1]$: any real value in the interval is possible.
\end{itemize}
\end{frame}

\begin{frame}{Two Types of Random Variables}
\begin{columns}[T,onlytextwidth]
\column{0.48\textwidth}
\begin{block}{Discrete random variable}
Takes values from a finite or countable set.

\medskip
Example:
\[
X\in\{0,1,2\}.
\]
\end{block}

\column{0.48\textwidth}
\begin{block}{Continuous random variable}
Can take any value in an interval.

\medskip
Example:
\[
0\le Y\le 1.
\]
\end{block}
\end{columns}

\vspace{0.35cm}
A measured variable can be intrinsically continuous even when recorded in rounded units.
\end{frame}

\begin{frame}{Parameters vs. Statistics}
\begin{columns}[T,onlytextwidth]
\column{0.48\textwidth}
\begin{block}{Sample statistics}
Computed from observed sample data:
\[
\bar{x},\qquad s.
\]
\end{block}

\column{0.48\textwidth}
\begin{block}{Distribution parameters}
Describe the theoretical probability distribution:
\[
\mu,\qquad \sigma.
\]
\end{block}
\end{columns}

\vspace{0.4cm}
The mean of a random variable is also called its \textbf{expected value}.
\end{frame}

\begin{frame}{Mean of a Discrete Random Variable}
Let $X$ have possible values $x_1,\ldots,x_k$ with probabilities
$p_1,\ldots,p_k$.

\begin{block}{Expected value}
\[
\mu=\E(X)=\sum_{i=1}^{k}x_i p_i.
\]
\end{block}

\begin{itemize}
  \item The expected value is a \textbf{probability-weighted average}.
  \item Outcomes with larger probabilities contribute more strongly to $\mu$.
  \item The value of $\mu$ does not need to be one of the possible individual outcomes.
\end{itemize}

\begin{block}{Long-run interpretation}
If the random experiment were repeated many times, the average observed value
of $X$ would tend toward $\mu$.
\end{block}
\end{frame}

\begin{frame}{Variance and Standard Deviation of a Random Variable}
Once the mean $\mu$ is known, we can describe how much the random variable
typically varies around that long-run average.

\begin{align*}
\sigma^2&=\sum_{i=1}^{k}(x_i-\mu)^2p_i,\\[0.35cm]
\sigma&=\sqrt{\sigma^2}.
\end{align*}

\begin{itemize}
  \item Each squared distance $(x_i-\mu)^2$ is weighted by its probability.
  \item The variance $\sigma^2$ is measured in squared units.
  \item Taking the square root gives $\sigma$, which is expressed in the same units as $X$.
\end{itemize}

\begin{alertblock}{Connection to sample statistics}
This is the theoretical analogue of the sample standard deviation $s$, but
here the calculation uses the entire probability distribution.
\end{alertblock}
\end{frame}

\begin{frame}{Example 9.13: Expected Hearing Impairment}
For $X\in\{0,1,2\}$ with probabilities $0.70,0.22,0.08$,
\begin{align*}
\mu
&=(0)(0.70)+(1)(0.22)+(2)(0.08)\\
&=\boxed{0.38}.
\end{align*}

\begin{alertblock}{Interpretation}
The expected value need not be a possible individual outcome. A dog cannot have $0.38$ impaired ears, but $0.38$ is the long-run average across many dogs.
\end{alertblock}
\end{frame}

\begin{frame}{Example 9.13: Variance and Standard Deviation}
\small
\begin{align*}
\sigma^2
&=(0-0.38)^2(0.70)+(1-0.38)^2(0.22)\\
&\quad +(2-0.38)^2(0.08)\\
&=0.10108+0.084568+0.209952\\
&=\boxed{0.3956},\\[0.2cm]
\sigma&=\sqrt{0.3956}=\boxed{0.62897}\approx 0.63.
\end{align*}
\end{frame}

\begin{frame}[fragile]{Example 9.13: Computing $\mu$ and $\sigma$ in R}
The formulas translate directly into vectorized R operations.

\vspace{0.15cm}
\scriptsize
\begin{verbatim}
# Define the outcomes (x) and their probabilities (p_x)
x <- c(0, 1, 2)
p_x <- c(0.70, 0.22, 0.08)

# 1. Calculate the Expected Value (mu)
mu <- sum(x * p_x)
print(mu)
# [1] 0.38

# 2. Calculate the Variance
variance <- sum((x - mu)^2 * p_x)

# 3. Calculate the Standard Deviation
sigma <- sqrt(variance)
print(sigma)
# [1] 0.6289674
\end{verbatim}

\normalsize
R performs the multiplication element-by-element, so \texttt{x * p\_x}
corresponds directly to the terms $x_iP(X=x_i)$ in the expected-value formula.
\end{frame}

% =============================================================
% PART VI: RISK AND ODDS
% =============================================================

\begin{frame}{Probability, Risk, and Odds}
In the life sciences, the same uncertain event may be communicated using
\textbf{probability}, \textbf{risk}, or \textbf{odds}. These quantities are
related, but they are not numerically identical.

If an event $A$ occurs with probability $p$, then:

\begin{block}{Risk}
\[
\operatorname{Risk}(A)=p.
\]
In health applications, risk usually refers to the probability of an
undesirable outcome.
\end{block}

\begin{block}{Odds}
\[
\operatorname{Odds}(A)=\frac{p}{1-p}.
\]
Odds compare the probability that the event occurs with the probability that
it does \emph{not} occur.
\end{block}
\end{frame}

\begin{frame}{Example 9.14: DVT as a Rare Event}
\small
Patients who remain immobilized for long periods can develop deep vein
thrombosis (DVT). In the study described in the notes, $42$ of $1{,}518$
patients receiving Fragmin experienced DVT.

\begin{align*}
p&=\frac{42}{1518}=0.0277 \qquad (2.77\%),\\[0.15cm]
\operatorname{Odds}
&=\frac{0.0277}{1-0.0277}
 =\frac{42}{1476}=0.0285.
\end{align*}

The odds can also be written approximately as
\[
1:35.
\]

\begin{block}{How to read $1:35$}
For approximately every one patient who experiences DVT, about $35$ do not.
\end{block}

\begin{alertblock}{Rare-event rule of thumb}
When $p$ is small, risk and odds are numerically similar because $1-p\approx 1$.
\end{alertblock}
\end{frame}

\begin{frame}{Example 9.15: Sickle-Cell Anemia}
If both parents carry the sickle-cell trait, the probability described in the
notes that a child inherits the disease is
\[
p=0.25.
\]

The \textbf{risk} is therefore $25\%$. The corresponding odds are
\begin{align*}
\operatorname{Odds}
&=\frac{0.25}{1-0.25}\\
&=\frac{0.25}{0.75}\\
&=\frac13\approx 0.333.
\end{align*}

Thus the odds are $1:3$: one occurrence for every three non-occurrences.

\begin{alertblock}{Important}
For more common events, risk and odds can differ substantially. Always check
which quantity a medical or scientific report is presenting.
\end{alertblock}
\end{frame}

\begin{frame}{Converting Odds Back to Probability}
Starting from
\[
\text{odds}=\frac{p}{1-p},
\]
solve for $p$:
\[
\boxed{p=\frac{\text{odds}}{1+\text{odds}}}.
\]

For odds of $1/3$,
\begin{align*}
p
&=\frac{1/3}{1+1/3}\\
&=\frac{1/3}{4/3}\\
&=\frac14=0.25.
\end{align*}
\end{frame}

% =============================================================
% PART VII: CHAPTER REVIEW
% =============================================================

\begin{frame}[fragile]{Summary of R Commands (Chapter 9)}
The chapter uses a small set of R tools to connect probability formulas with
simulation and computation.

\begin{table}
\centering
\scriptsize
\begin{tabular}{lp{9.1cm}}
\toprule
\textbf{R command} & \textbf{Role in this chapter} \\
\midrule
\texttt{sample(..., replace = TRUE)} & Simulates repeated random outcomes, such as coin tosses. \\
\addlinespace
\texttt{cumsum(x)} & Creates a running total, useful for visualizing long-run relative frequency. \\
\addlinespace
\texttt{data.frame(...)} & Stores outcomes and their probabilities together in a probability-model table. \\
\addlinespace
\texttt{sum(x * p\_x)} & Computes the expected value of a discrete random variable. \\
\addlinespace
\texttt{sum((x - mu)\string^2 * p\_x)} & Computes the probability-weighted variance. \\
\addlinespace
\texttt{sqrt(variance)} & Converts variance into standard deviation. \\
\addlinespace
\texttt{ggplot()}, \texttt{geom\_bar()}, \texttt{geom\_rect()} & Visualize discrete probabilities and continuous probability areas. \\
\bottomrule
\end{tabular}
\end{table}
\end{frame}

\begin{frame}{Key Takeaways}
\begin{itemize}
  \item Probability describes uncertainty using numbers from $0$ to $1$.
  \item Long-run relative frequency provides the frequentist interpretation.
  \item A probability model consists of a sample space and probability assignments.
  \item Disjoint events can be added; complements are found by subtracting from $1$.
  \item Discrete probabilities are sums; continuous probabilities are areas.
  \item Random variables have theoretical means and standard deviations.
  \item Risk is a probability, while odds are $p/(1-p)$.
\end{itemize}
\end{frame}

\begin{frame}{Check Your Understanding}
\begin{enumerate}
  \item If $\prob{A}=0.37$, what is $\prob{A^c}$?
  \item If $A$ and $B$ are disjoint with probabilities $0.25$ and $0.40$, what is $\prob{A\text{ or }B}$?
  \item For a continuous random variable, what is $\prob{X=5}$?
  \item If $p=0.20$, what are the odds of the event?
\end{enumerate}

\pause
\vspace{0.4cm}
\begin{block}{Answers}
$0.63$; $0.65$; $0$; $0.20/0.80=0.25$, or $1:4$.
\end{block}
\end{frame}


\end{document}
